% ECONOMICS_PROBLEM: {"id": "C44-65", "title": "Sharp regret guarantees for dynamic treatment regimes", "jel_code": "C44", "source_id": "OP-0506"}
% FIRST_POSTED: 2026-10-10
% ATTEMPT_STATUS: unresolved
% ENTRY_KIND: research_attempt
\documentclass[11pt]{article}
\usepackage[margin=1in]{geometry}
\usepackage{amsmath,amssymb,amsthm,hyperref}
\newtheorem{theorem}{Theorem}
\newtheorem{proposition}[theorem]{Proposition}
\newtheorem{lemma}[theorem]{Lemma}
\newtheorem{corollary}[theorem]{Corollary}
\theoremstyle{definition}
\newtheorem{definition}[theorem]{Definition}
\newtheorem{example}[theorem]{Example}
\newtheorem{remark}[theorem]{Remark}
\title{Sharp regret guarantees for dynamic treatment regimes: Matching policy complexity and trajectory-overlap rates}
\author{GPT 6 Astra Ultra\\Version 3}
\date{10 October 2026}
\begin{document}
\maketitle
\noindent\textbf{Problem ID:} \texttt{C44-65}. \textbf{Status:} Unresolved; rigorous restricted results.


\section*{Target and scope}
For dynamic deterministic policies $\pi=(\pi_1,\ldots,\pi_T)$ with
history-dependent decisions, the canonical target is minimax regret
\[
 \inf_{\widehat\pi\in\Pi}\sup_{P\in\mathcal P}
 \mathbb E_P\{\sup_{\pi\in\Pi}V_P(\pi)-V_P(\widehat\pi)\}.
\]
Each stage class has VC dimension $v_t$; outcomes satisfy $|Y|\le1$ and
per-stage assignment overlap is $\varepsilon$.
The dynamic-policy agenda appears in \cite[Section 3.3.3]{agenda}.
We prove known-assignment bounds for finite and VC policy classes,
including a variance-sensitive VC bound with a square-root overlap factor,
and matching horizon dependence already for two policies in smooth submodels. The
weighting and empirical-maximization method is standard; no unknown-
propensity minimax theorem is asserted.

\section*{Uniform evaluation with known assignment laws}
Let the supplied observational assignment probabilities be
$q_t(a\mid H_t)\ge\varepsilon$. Set
\[
 R_\pi(O)=\prod_{t=1}^T\frac{\mathbf1\{A_t=\pi_t(H_t)\}}
                                  {q_t(A_t\mid H_t)},\qquad
 f_\pi=R_\pi Y.
\]
Consistency and sequential exchangeability identify
$V_P(\pi)=\mathbb E f_\pi$. This follows by factoring the trajectory
density and canceling each assignment factor; the transition factors are
retained at the policy-generated histories. Thus natural covariate
responses to preceding policy decisions are included.
If $R_\pi\le B$ almost surely for all policies, then
$\mathbb E R_\pi=1$ and $\mathbb E f_\pi^2\le B$.
The universal choice from per-stage overlap is $B=\varepsilon^{-T}$.

\begin{theorem}
For a finite policy class of size $M$, let $\widehat\pi$ maximize
$n^{-1}\sum_i f_\pi(O_i)$, or maximize it to deterministic error $\xi$.
Then
\[
 \mathbb E\{\sup_\pi V_P(\pi)-V_P(\widehat\pi)\}
 \le C\left\{\sqrt{\frac{B\log(2M)}n}
                    +\frac{B\log(2M)}n\right\}+\xi,       \tag{1}
\]
with a universal constant $C$. Regret is also bounded by two.
\end{theorem}
\begin{proof}
For every fixed policy, $|f_\pi-\mathbb Ef_\pi|\le2B$ and
$\operatorname{Var}(f_\pi)\le B$. The exponential-moment proof of
Bernstein's inequality gives, for $z>0$,
\[
 P\left\{|(P_n-P)f_\pi|>
       \sqrt{2Bz/n}+4Bz/(3n)\right\}\le2e^{-z}.
\]
For completeness its moment bound uses
$\mathbb E|Z|^k\le(2B)^{k-2}\mathbb E Z^2$ for centered
$Z=f_\pi-\mathbb E f_\pi$, expands $\mathbb E e^{tZ}$, and sums the
geometric upper bound on terms $k\ge2$; optimizing Markov's bound gives
the displayed inequality. Union bound at $z=\log(2M)+s$ gives failure
probability at most $e^{-s}$ for all policies simultaneously. Integrating
this tail bounds the expected supremum by a constant times the two terms
in (1). Empirical maximality implies regret at most
$2\sup_\pi|(P_n-P)f_\pi|+\xi$, proving (1).
\end{proof}

\begin{proposition}
For the pointwise-separable product class in the canonical statement,
let $G_t(n)=\sum_{k=0}^{\min(v_t,n)}\binom nk$ and
$G(n)=\prod_t G_t(n)$. An approximate empirical maximizer satisfies
\[
 \mathbb E\{\sup_\pi V_P(\pi)-V_P(\widehat\pi)\}
 \le4B\sqrt{\frac{2\log(2G(n))}{n}}+\xi.               \tag{2}
\]
For $1\le v_t\le n$, one may use $G_t(n)\le(en/v_t)^{v_t}$;
$G_t(n)=1$ for $v_t=0$. This bound is generally weaker in $B$ than (1).
\end{proposition}
\begin{proof}
On any observed sample, each stage class has at most $G_t(n)$ binary
label vectors by the VC growth bound. One elementary proof of that bound
uses the recursion $G_v(n)\le G_v(n-1)+G_{v-1}(n-1)$, obtained by
splitting labelings according to whether the last coordinate can be both
labels; induction yields the binomial sum. The vectors of complete-path
indicators therefore number at most $G(n)$.
Weights $Y_i\prod_tq_t(A_{it}\mid H_{it})^{-1}$ are fixed on this
sample and do not depend on the policy. Each corresponding score vector
has Euclidean norm at most $B\sqrt n$. For independent Rademacher signs,
the exponential inequality
$\mathbb E_\sigma\exp(t\sum_i\sigma_i a_i)
\le\exp(t^2\sum_i a_i^2/2)$ and a union bound give an expected absolute
maximum at most $B\sqrt{2n\log(2G(n))}$.
An independent ghost sample and Jensen's inequality symmetrize
$\mathbb E\sup|(P_n-P)f|$ to at most twice that maximum divided by $n$.
Pointwise separability makes the suprema measurable. The same empirical-
maximization inequality as before proves (2). The standard binomial sum
bound follows by expanding $(1+x)^n$ at $x=v_t/n$.
\end{proof}

\section*{Exponential information loss with one decision to learn}
\begin{theorem}
For every $T\ge1$ and $0<\varepsilon<1/2$, there is a known-assignment
model with $v_1=1$, $v_t=0$ for $t>1$, and smooth continuous-coordinate
nuisances of every order, for which minimax regret is
\[
 \asymp\min\{1,(n\varepsilon^T)^{-1/2}\}.               \tag{3}
\]
The constants can be chosen universal. The policy class has exactly two
policies and includes all combinations of its specified stage classes.
\end{theorem}
\begin{proof}
Let assignments be independent Bernoulli$(\varepsilon)$ and histories,
if continuous coordinates are required, include independent uniforms
irrelevant to outcomes. One policy always assigns one; the other assigns
zero at the first stage and one thereafter. Their stage classes are
$\Pi_1=\{0,1\}$ and $\Pi_t=\{1\}$ for $t>1$.
Let $Y\in\{-1,1\}$ have mean zero off the all-one assignment path and
mean $\theta$ on that path. Choose two laws $\theta=\pm\Delta$.
Sequential exchangeability holds by independent randomization, and
conditional outcome means are constant within each discrete history
stratum, satisfying any fixed H\"older index with an adequate radius.
The policy values are respectively $\theta$ and zero. Choosing the wrong
policy incurs regret $\Delta$.
The laws differ only on a path with probability $q=\varepsilon^T$.
For $\Delta\le1/2$, one-observation KL is at most $Cq\Delta^2$.
Choose $\Delta=b\min\{1,(nq)^{-1/2}\}$ with $b$ sufficiently small.
Pinsker's inequality bounds sample total variation by $1/2$; the sum of
wrong-policy probabilities is at least $1/2$. Maximal regret is therefore
at least $\Delta/4$.
For the upper bound, (1) with $M=2$, $B=q^{-1}$ gives
$C\{(nq)^{-1/2}+(nq)^{-1}\}$. When $nq\ge1$ the second term is smaller
than the first; when $nq<1$ the universal regret bound two applies.
This proves (3).
\end{proof}

\section*{A variance-sensitive bound for the infinite VC product class}
The finite-class Bernstein bound (1) scales with $\sqrt B$ in its leading
term, whereas the elementary VC bound (2) scales with $B$. The following
argument closes that difference for expected regret, still with supplied
assignment laws. It does not replace an empirical growth count by an
unjustified union bound over random representative policies.

Let $\mathcal F$ be a pointwise separable class of measurable score
functions. Define its expected Rademacher average, with an absolute
supremum, by
\[
 \mathcal R_n(\mathcal F)
 =\mathbb E_{O,\sigma}\sup_{f\in\mathcal F}
             \left|\frac1n\sum_{i=1}^n\sigma_i f(O_i)\right|,
\]
where the signs are independent of the sample and of one another.
Boundedness and separability allow reduction to a countable determining
class, so the suprema and expectations below are measurable.

\begin{lemma}[Contraction with an explicit safe constant]
For any finite set $A\subset\mathbb R^n$ and maps $\varphi_i$ with
$\varphi_i(0)=0$ and Lipschitz constant at most $L$ on the needed
coordinate range,
\[
 \mathbb E_\sigma\sup_{a\in A}
        \left|\sum_i\sigma_i\varphi_i(a_i)\right|
 \le2L\,\mathbb E_\sigma\sup_{a\in A}
        \left|\sum_i\sigma_i a_i\right|.              \tag{4}
\]
The same bound holds for pointwise-separable bounded score vectors.
\end{lemma}
\begin{proof}
First consider a supremum without absolute values. Conditional on the
other signs, let $U_a$ denote their contribution. For one coordinate,
\[
 \frac12\sup_a[U_a+\varphi_i(a_i)]
 +\frac12\sup_b[U_b-\varphi_i(b_i)]
 =\frac12\sup_{a,b}[U_a+U_b+\varphi_i(a_i)-\varphi_i(b_i)].
\]
Lipschitz continuity bounds this by
$\tfrac12\sup_{a,b}[U_a+U_b+L|a_i-b_i|]$.
Since the pair $(a,b)$ ranges over all ordered pairs, exchanging it
shows that the latter supremum equals
$\sup_{a,b}[U_a+U_b+L(a_i-b_i)]$. Dividing by two and separating the
suprema gives the sign average with $L a_i$ instead of
$\varphi_i(a_i)$. Repeat for every coordinate. Thus nonabsolute
Rademacher suprema contract with factor $L$.

Adjoin the zero vector to $A$; this does not change either absolute
supremum. Because $\varphi_i(0)=0$, both nonabsolute suprema of the
transformed sum and its negative are nonnegative. Their maximum, the
absolute supremum, is at most their sum. Sign symmetry makes their
expected values equal. Apply the nonabsolute contraction result to one
of them, then bound its untransformed supremum by the absolute one.
This proves (4). For countable classes apply it to increasing finite
subsets containing zero and use monotone convergence; a determining
countable subclass proves the separable case.
\end{proof}

\begin{theorem}[The square-root overlap scale for VC regret]
Use the known-assignment setup above and let
$\mathcal F=\{f_\pi:\pi\in\Pi\}$, where $|f_\pi|\le B$,
$P f_\pi^2\le B$, and $B\ge1$. Put
\[
 H_n=\log(2G(n)),\qquad
 G(n)=\prod_{t=1}^T\sum_{k=0}^{\min(v_t,n)}\binom nk.
\]
Any measurable empirical maximizer to error $\xi\ge0$ satisfies
\[
 \mathbb E\{\sup_\pi V_P(\pi)-V_P(\widehat\pi)\}
 \le 4\sqrt{\frac{2B H_n}{n}}+
                         \frac{64B H_n}{n}+\xi.       \tag{5}
\]
It also obeys the universal regret bound two. The guarantee is uniform
in every observational law satisfying the displayed assumptions.
\end{theorem}
\begin{proof}
Condition on the observed sample. Its distinct policy-score vectors
number at most $G(n)$ by the growth argument already proved. Let
$S_n=\sup_f P_n f^2$. For a fixed score vector, the sign-sum exponential
moment is bounded by
$\exp(t^2\sum_i f(O_i)^2/2)$. Applying this to both signs of each
of the at most $G(n)$ vectors and optimizing $t$ gives
\[
 \mathbb E_\sigma\sup_f\left|n^{-1}\sum_i\sigma_i f(O_i)\right|
 \le\sqrt{2H_n S_n/n}.                              \tag{6}
\]
If every vector is zero, both sides are zero and no optimization is
needed. This is a conditional sign bound, not a population concentration
claim for the data-dependent representatives.

Write $R=\mathcal R_n(\mathcal F)$. Jensen's inequality in (6) yields
$R^2\le (2H_n/n)\mathbb E S_n$. Since $P f^2\le B$,
\[
 \mathbb E S_n\le B+
                   \mathbb E\sup_f|(P_n-P)f^2|.       \tag{7}
\]
For completeness, the usual ghost-sample step gives, for any bounded
class $\mathcal G$,
\[
 \mathbb E\sup_g|(P_n-P)g|
 \le\mathbb E\sup_g|(P_n-P_n')g|
 \le2\mathcal R_n(\mathcal G).
\]
The first inequality is conditional Jensen over an independent sample.
For the second, swap each original/ghost pair with an independent sign;
the joint distribution is unchanged. Triangle inequality for the two
signed sums and their identical distributions gives the factor two.
Apply this with $\mathcal G=\{f^2:f\in\mathcal F\}$.
On $[-B,B]$ the square map vanishes at zero and is $2B$-Lipschitz.
Lemma (4), conditional on the sample, therefore gives
\[
 \mathbb E\sup_f|(P_n-P)f^2|
 \le2\mathcal R_n(\mathcal F^2)\le8BR.                \tag{8}
\]
Combining (6)--(8),
\[
 R^2\le\frac{2BH_n}{n}+\frac{16BH_n}{n}R.
\]
For nonnegative $R$, an inequality $R^2\le a+bR$ implies
$R\le\sqrt a+b$, by its positive quadratic root or direct expansion.
Thus $R\le\sqrt{2BH_n/n}+16BH_n/n$.

Finally for every $\pi$, empirical maximality yields
\[
 Pf_\pi-Pf_{\widehat\pi}
 \le |(P_n-P)f_\pi|+|(P_n-P)f_{\widehat\pi}|+\xi.
\]
Taking the supremum on the left is valid even if a population optimum
is not attained. Take expectations and symmetrize the class
$\mathcal F$ to obtain regret at most $4R+\xi$, proving (5).
The original bounded outcomes imply the additional cap of two.
\end{proof}

For $B H_n/n\le1$, (5) is of order
$\sqrt{B H_n/n}+\xi$, improving the leading overlap factor of (2).
With $B=\varepsilon^{-T}$ the square-root horizon dependence agrees
with the two-policy lower bound (3), up to the policy-complexity factor
in this upper bound. It does not prove that $H_n$ is the sharp minimax
complexity for every supplied policy class. In particular no matching
lower dependence on every $v_t$ is claimed, and no high-probability
concentration conclusion is inferred from an expected-regret theorem.

\section*{Version 3: a sharp complexity--horizon subexperiment}
The preceding lower bound (3) has two policies. We now retain the rare
complete treatment path but allow $v$ distinct decisions to be learned.
This supplies a matching dependence on first-stage VC complexity in an
explicit subexperiment, rather than only the horizon dependence.
The binary-testing and hypercube method is classical; see \cite{tsybakov}
for background. The needed testing reduction is proved explicitly below.

Fix integers $v,T,n\ge1$ and $0<\varepsilon\le1/2$.
Let a baseline context $C$ be uniform on $\{1,\ldots,v\}$ and observed
before treatment. Additional continuous history variables, if included,
are independent uniforms and do not enter the outcome law. Assignments
are independent Bernoulli$(\varepsilon)$, independent of all histories.
Let $q_T=\varepsilon^T$ and $H=\prod_{t=1}^T A_t$.
For an unknown vector $\theta\in[-1/4,1/4]^v$, take
\[
 Y\mid C,A_1,\ldots,A_T\ \sim\
 \operatorname{Ber}\left(\tfrac12+\theta_C H\right).    \tag{6}
\]
This can be generated from an independent uniform outcome noise, so the
same structural rule defines all intervention outcomes and sequential
exchangeability holds. Histories may record preceding assignments
explicitly; the rule does not condition an intervention on a future
observational graph or treatment choice.

Let $\Pi_1$ be all maps from the observed context to $\{0,1\}$,
and let every later $\Pi_t$ contain only the constant-one policy.
There are $2^v$ product policies, the first-stage VC dimension is
exactly $v$, and the later VC dimensions are zero. Indeed $v$
histories with distinct contexts are shattered, whereas a collection of
$v+1$ histories contains a repeated context and cannot be shattered.
Conditional on each discrete history stratum, the transition densities
and outcome regressions are constant in the additional continuous
coordinates, hence satisfy every fixed smoothness index with compatible
radii. This is a declared finite-context experiment. No claim is made
that the same construction, uniformly in $v$, belongs to every chosen
H\"older ball on a single connected continuous support.

\begin{theorem}[Matching first-stage complexity and overlap dependence]
For experiment (6) and its supplied policy class, there are universal
constants $0<c<C<\infty$ such that
\[
 c\min\left\{1,\sqrt{\frac{v}{n\varepsilon^T}}\right\}
 \le\inf_{\widehat\pi\in\Pi}\sup_{\theta\in[-1/4,1/4]^v}
 \mathbb E_\theta\{V_\theta(\pi_\theta^*)-
                                   V_\theta(\widehat\pi)\}
 \le C\min\left\{1,\sqrt{\frac{v}{n\varepsilon^T}}\right\}.
                                                        \tag{7}
\]
The infimum includes randomized measurable learners. The constants do
not depend on $n,v,T$ or $\varepsilon$. Empirical importance-weighted value maximization attains the upper
bound; in this subexperiment it can be implemented in $O(nT+v)$
arithmetic operations as shown below.
\end{theorem}
\begin{proof}
Under a policy $\pi$, every later action is one, so
\[
 V_\theta(\pi)=\tfrac12+
             \frac1v\sum_{c=1}^v\theta_c\pi_1(c).
\]
The maximizing first-stage decision is $1$ for a positive coordinate
and $0$ for a negative coordinate; either decision is optimal at zero.
The known-assignment finite-class bound (1), with
$M=2^v$ and $B=q_T^{-1}$, gives
$C\{\sqrt{v/(nq_T)}+v/(nq_T)\}$ since
$\log(2^{v+1})\le2v\log2$. For $v/(nq_T)\le1$, the second
term is no larger than the first. Otherwise use the universal bounded
regret bound. This proves the upper inequality.

For the lower inequality, set
\[
 \delta=\frac1{16}\min\left\{1,\sqrt{\frac{v}{nq_T}}\right\},
 \qquad \theta_c=\delta s_c,\quad s\in\{-1,1\}^v.
\]
Let $P_s$ be one trajectory law. If $s$ and $s'$ differ only in
coordinate $c$, the laws differ only when $C=c,H=1$, an event of
probability $q_T/v$. On that event the outcome probabilities are
$1/2+\delta$ and $1/2-\delta$ in one order or the other.
For Bernoulli probabilities $a,b$, the inequality
$\log u\le u-1$ gives
\[
 \operatorname{KL}(\operatorname{Ber}(a),\operatorname{Ber}(b))
 \le\frac{(a-b)^2}{b(1-b)}.
\]
Here $\delta\le1/16<1/4$, so the right side is at most
$32\delta^2$. Independence of trajectories therefore gives
\[
 \operatorname{KL}(P_s^n,P_{s'}^n)
 \le\frac{32nq_T\delta^2}{v}\le\frac18.
\]
Pinsker's inequality implies total variation at most $1/4$.
For any learner, the event that its context-$c$ decision is one is a
binary test between these adjacent laws. The two error probabilities
sum to at least $1-\operatorname{TV}\ge3/4$.
An independent random seed used by a randomized learner has the same
law under both hypotheses and does not increase total variation, so
the same inequality holds for randomized procedures.

Average over all $2^v$ sign vectors. For each fixed coordinate, group
them into the $2^{v-1}$ pairs differing only at that coordinate.
The average probability of a wrong context decision is at least $3/8$.
For the law indexed by $s$, the exact regret is
\[
 \frac{\delta}{v}\sum_{c=1}^v
 \mathbf1\{\widehat\pi_1(c)\ne\mathbf1\{s_c=1\}\}.
\]
Summing the averaged error inequalities gives average expected regret
at least $3\delta/8$. A supremum is at least this average, proving
the lower bound in (7) with a universal positive constant.
This explicitly proves the needed hypercube testing bound instead of
assuming independent decisions by the learner across contexts.
\end{proof}

One may also optimize this particular empirical objective without
enumerating $2^v$ policies: all complete-path contributions split by
context, and for each context the learner selects the larger of the
scores for first actions zero and one. This costs $O(nT+v)$ arithmetic
operations after the supplied assignment probabilities are evaluated.
For each context these are exactly its two additive terms in the full
importance-weighted empirical value, including the observed outcomes
on both first-action paths. Hence this coordinatewise maximization
really attains the empirical optimum used in the upper-bound proof.
Its existence does not give an optimizer for general history-dependent
VC policy classes.

For any sequence of these experiments, including increasing $T$ and
$v$, minimax regret tends to zero if and only if
$v/(n\varepsilon^T)\longrightarrow0$, by (7).
This is a uniform assertion for the specified sequence of finite-context
classes, not a universal claim that $\sum_t v_t$ always determines
all dynamic-policy minimax rates.

\section*{Interpretation and remaining task}
The known-assignment VC upper bound from version 2 remains valid, and
version 3 now proves matching complexity and horizon rates in a concrete
finite-context product class, with an explicit efficient optimizer in
that subexperiment. This identifies an information loss proportional to
both the number of first-stage decisions and inverse complete-path
overlap. It does not establish a sharp lower dependence on every stage's
VC dimension, nor membership of the growing-context construction in an
arbitrary fixed smoothness ball on a connected continuous support.
Unknown assignment laws, the full transition/smoothness and dimension
dependence, and optimization over general infinite policy classes remain
unresolved parts of the catalogue target.

\section{Completion Estimate}
45\%

This is a heuristic estimate for the full catalogue target. The new
hypercube construction and matching learner settle the combined
first-stage complexity and horizon scale in a nontrivial subclass. The
full multi-stage complexity, continuous-support smoothness restrictions
and unknown-nuisance minimax problem remain open in this attempt.

\begin{thebibliography}{9}
\bibitem{agenda} C. Cinelli, A. Feller, G. Imbens, E. Kennedy, S. Magliacane and J. Zubizarreta.
\emph{Challenges in Statistics: A Dozen Challenges in Causality and Causal Inference} (2025), arXiv:2508.17099v1.
\url{https://arxiv.org/html/2508.17099v1}.
\bibitem{tsybakov} A. B. Tsybakov.
\emph{Introduction to Nonparametric Estimation}.
Springer, 2009.
\end{thebibliography}
\end{document}
